\input{../tex-inputs/latex-leseansicht-vorspann}

\section[Gustav Schwarzkopf an Arthur Schnitzler, 31. 8. 1899]{L04294 Gustav Schwarzkopf an Arthur Schnitzler, 31. 8. 1899}
\nopagebreak\mylabel{L04294v}
\rehead{ }\normalsize\beginnumbering\briefempfaengerindex{Schnitzler, Arthur@\textsc{Schnitzler, Arthur}!zzzSchwarzkopf, Gustav@\emph{von Gustav Schwarzkopf}!1899-08-312@{31. 8. 1899}|(be}
\toendnotes[C]{\smallbreak\pagebreak[2]}
\correspDesc{Versand  durch Gustav Schwarzkopf am 31. 8. 1899 in Wien
\newline{}Erhalt  durch Arthur Schnitzler im Zeitraum [1. 9. 1899 – 5. 9. 1899?] in Bad Ischl}\toendnotes[C]{\smallbreak}
\Standort{CUL, Schnitzler, B 96a.}
\physDesc{Briefkarte, 1135 Zeichen
\newline{}Handschrift: schwarze Tinte, deutsche Kurrent}\toendnotes[C]{\smallbreak}
\pstart
           \raggedleft{}{\pb}31. 8. 99.\pend
           \vspace{0.5em}
\pstart
           Lieber Arthur, ich will Ihnen nur anzeigen, daß ich noch lebe und
               Ihnen für Ihren lieben \label{K_L04294-1v}\edtext{Brief}{\lemma{\textnormal{\emph{Brief}}}\Cendnote{\textnormal{XXXX Auszeichnungsfehler: Dokument L04131 nicht gefunden. }}}\label{K_L04294-1} danken. Ich hatte die
               redliche Abſicht Ihnen auch{ }ſechs Seiten zu{ }ſchreiben, die notwendige Zeit hätte ich{ }ſchon dazu erübrigt, aber es fehlt mir wirklich an Material. Nichts, das eine
               Mitteilung forderte, nichts, das auch nur eine Betrachtung verdient. In den letzten
               vier Wochen habe ich mich nicht einmal mehr geärgert, nicht einmal über die
               geſchmacklose Art die N. P.\pwindex{Neue Freie Presse@\emph{Neue Freie Presse}|pw} in der Dreyfus\pwindex{Dreyfus, Alfred 9.\,10.\,1859 Mulhouse – 12.\,7.\,1935 Paris@\textsc{Dreyfus, Alfred} (9.\,10.\,1859 Mulhouse – 12.\,7.\,1935 Paris), \emph{Militär}|pw}-Sache, {\pb}die einem{ }ſelbſt das Mitleid verleiden
                  kö{\geminationn}te, nicht einmal über die \label{K_L04294-2v}\edtext{Göthe\pwindex{Goethe, Johann Wolfgang von 28.\,8.\,1749 Frankfurt am Main – 22.\,3.\,1832 Weimar@\textsc{Goethe, Johann Wolfgang von} (28.\,8.\,1749 Frankfurt am Main – 22.\,3.\,1832 Weimar), \emph{Schriftsteller}|pw}-Feier}{\lemma{\textnormal{\emph{Göthe-Feier}}}\Cendnote{\textnormal{Am
                     28. 8. 1899 hätte Goethe\pwindex{Goethe, Johann Wolfgang von 28.\,8.\,1749 Frankfurt am Main – 22.\,3.\,1832 Weimar@\textsc{Goethe, Johann Wolfgang von} (28.\,8.\,1749 Frankfurt am Main – 22.\,3.\,1832 Weimar), \emph{Schriftsteller}|pwk} seinen 150. Geburtstag gehabt. Im Unterschied zu anderen Bühnen,
                  die zeitnah Festakte inszenierten, wollte das \emph{Burgtheater}\orgindex{Burgtheater@Burgtheater|pwk} einen solchen erst am 8. 10. 1899\eventindex{Burgtheater@\textbf{Burgtheater}!Goethe-Feier, 8.10.1899@Goethe-Feier, 8.10.1899|pwkv} veranstalten. Nachdem
                  Kritik laut wurde, wurde die Saisoneröffnung mit \emph{Götz von
                        Berlichingen}\pwindex{Götz von Berlichingen mit der eisernen Hand@\emph{Götz von Berlichingen mit der eisernen Hand}|pwk} am 1. 9. 1899\eventindex{Burgtheater@\textbf{Burgtheater}!Aufführung von Götz von Berlichingen mit der eisernen Hand, 1.9.1899@Aufführung von Götz von Berlichingen mit der eisernen Hand, 1.9.1899|pwkv} zur Göthe\pwindex{Goethe, Johann Wolfgang von 28.\,8.\,1749 Frankfurt am Main – 22.\,3.\,1832 Weimar@\textsc{Goethe, Johann Wolfgang von} (28.\,8.\,1749 Frankfurt am Main – 22.\,3.\,1832 Weimar), \emph{Schriftsteller}|pwk}-Feier ›umgedeutet‹. }}}\label{K_L04294-2} und die
               verſchiedenen Aufſätze. Es{ }ſcheint, daß auch zum Ärgern Geſellſchaft gehört. – Daß
               ich schon zweimal im Theater war, bei einer{ }ſehr{ }ſchlechten Vorſtellung\eventindex{Volkstheater@\textbf{Volkstheater}!Aufführung von Der Meineidbauer, 14.8.1899@Aufführung von Der Meineidbauer, 14.8.1899|pwv} (Meineidbauer\pwindex{\textcolor{red}{\textsuperscript{XXXX indx1}}!Meineidbauer. Volksstück mit Gesang in drei Acten@\strich\emph{Der Meineidbauer. Volksstück mit Gesang in drei Acten}|pw} und einer nicht guten\eventindex{Volkstheater@\textbf{Volkstheater}!Aufführung von Die Haubenlerche, 26.8.1899@Aufführung von Die Haubenlerche, 26.8.1899|pwv} (Haubenlerche\pwindex{\textcolor{red}{\textsuperscript{XXXX indx1}}!Haubenlerche@\strich\emph{Die Haubenlerche}|pw}) wird Sie kaum intereſſiren. Mich hat es auch nicht
               intereſſirt. Geſtern Abend habe ich in meiner Straße\oindex{Wien@\textbf{Wien}!I., Innere Stadt@\textbf{I., Innere Stadt}!Tiefer Graben@\textbf{Tiefer Graben}, \emph{Straße}|pwv} den D\textsuperscript{r}{ }Schickh\pwindex{Schik, Friedrich *~6.\,9.\,1857 Wien@\textsc{Schik, Friedrich} (*~6.\,9.\,1857 Wien), \emph{Notar, Journalist, Dramaturg}|pwu} getroffen, der jetzt
               wieder{ }ſehr »lieb« tut. Er läßt Sie grüßen. Ich{ }ſtelle mir vor, daß es vom
                  erſten September an{ }ſehr hübſch in \textsc{Ischl\oindex{Bad Ischl@\textbf{Bad Ischl}|pw}}{ }ſein wird; da müſſen doch{ }ſchon viele abgereiſt{ }ſein.\pend
           
\pstart
           Viel Vergnügen in Innsbruck\oindex{Innsbruck@\textbf{Innsbruck}, \emph{Verwaltungsgebiet}|pw}.\pend
           
\pstart
           Herzlichſt grüßt{\\[\baselineskip]}Ihr \spacefill\mbox{GustavS}\pend
           \leftskip=0em{}\selectlanguage{ngerman}\endnumbering\briefempfaengerindex{Schnitzler, Arthur@\textsc{Schnitzler, Arthur}!zzzSchwarzkopf, Gustav@\emph{von Gustav Schwarzkopf}!1899-08-312@{31. 8. 1899}|)be}\mylabel{L04294h}
\begin{anhang}
\end{anhang}\newcommand{\dateiname}{L04294}\newcommand{\titel}{Gustav Schwarzkopf an Arthur Schnitzler, 31. 8. 1899}\newcommand{\editorInnen}{Herausgegeben von Jahnke, SelmaMüller, Martin Anton}\input{../tex-inputs/latex-leseansicht-abspann}
